% ----------------------------------------------------------
% Introdução (exemplo de capítulo sem numeração, mas presente no Sumário)
% ----------------------------------------------------------

\chapter{Introdução}\label{intro}
A Análise do Comportamento Aplicada (ABA) é um método que se dedica à modelagem de comportamento, reforçando hábitos desejados e modificando aqueles que se apresentam problemáticos. Esse método envolve a utilização de atividades estimulantes de comportamentos, fornecendo diretrizes para promover a compreensão, bem como a observação do comportamento da criança após a sessão. Adicionalmente, reforçam-se ações que geraram atitudes positivas e incorporam-se intervalos de espera entre as atividades para proporcionar uma sessão mais tranquila.

No entanto, este tratamento é conduzido por um(a) terapeuta que realiza uma avaliação do comportamento e registra cuidadosamente suas observações sobre o paciente \cite{artoni2018technology}. Essas anotações são frequentemente feitas de forma manuscrita em papel, visando à consulta posterior. Durante o processo terapêutico, torna-se necessário efetuar diversas anotações, o que pode requerer um certo esforço, podendo, em alguns casos, resultar em três ou quatro folhas preenchidas \cite{artoni2013portable}.

Segundo uma pesquisa em 2020 da \textit{Centers for Disease Control and Prevention (CDC)} 1 a cada 36 crianças com 8 anos tem autismo \cite{fitch2023cdc}. Por conta disso, este trabalho oferece uma abordagem inovadora para o levantamento de dados aplicados ao método ABA, concentrando-se na eficiência e na automação de processos manuais. Essa abordagem visa aprimorar a precisão e reduzir erros. Além disso, busca melhorar a análise de dados, integrando anotações individuais e dados compilados para proporcionar uma visão global mais abrangente sobre o sucesso da aplicação do método ABA.

% Este trabalho, por outro lado, explora os métodos clássicos de otimização multiobjetivo, e está estruturado da seguinte forma: no Capítulo \ref{referencial_teorico} é apresentada uma breve explicação sobre o problema de otimização multiobjetivo, dos métodos clássicos e as principais famílias de teste, que se propõe a avaliar a qualidade dos algoritmos de otimização. No capítulo \ref{cap:metodologia} são apresentados os materiais utilizados e pseudocódigos dos métodos implementados e as particularidades de cada algoritmo. No capítulo \ref{cap:resultados} os resultados são discutidos, e por fim, têm-se as considerações finais no capítulo \ref{cap:conclusao} .

\section{Justificativa}
    
% A eficácia no tratamento de crianças com autismo não se resume apenas às intervenções de terapeutas especializados, mas exige a participação ativa da família \cite{maurice1996behavioral}. A literatura ressalta a contribuição significativa que os pais e cuidadores desempenham no progresso das crianças autistas, enfatizando a necessidade de uma abordagem colaborativa e sistemática.

A eficácia no tratamento de crianças com autismo depende da superação de diversos desafios enfrentados por terapeutas e familiares. Entre esses desafios, destacam-se a escassez de recursos humanos e tecnológicos que auxiliem na minimização de erros e na identificação da intervenção mais adequada para cada caso, o diagnóstico tardio, a falta de incentivos para pesquisas sobre o tema e, sobretudo, a ausência de um acompanhamento sistemático por parte da família. Como destacado por Maurice et al. (1996) \cite{maurice1996behavioral}, o tratamento do autismo não se restringe às intervenções realizadas por terapeutas especializados, mas exige a participação ativa da família. A literatura ressalta a importância dos pais e cuidadores no progresso das crianças autistas, reforçando a necessidade de uma abordagem colaborativa e sistemática.

% Fim correção

Além disso, referências, como o estudo de \cite{artoni2011accessible} sobre a viabilidade de softwares de comunicação para pessoas com dificuldades comportamentais ou de interação social, indicam a presença de soluções tecnológicas promissoras. No entanto, essa motivação vai além. Este estudo tem como objetivo desenvolver um aplicativo móvel que atuará como uma ferramenta eficaz para o gerenciamento de anotações por parte dos responsáveis no tratamento de crianças com autismo. O software proposto visa facilitar a comunicação e a troca de informações entre os cuidadores, simplificando a coleta de dados durante as sessões de terapia.


\section{Hipótese}
     Qual a viabilidade do desenvolvimento de softwares na adoção de metodologia \textit{Applied Behavior Analysis} (ABA) para auxiliar psicólogos no acompanhamento de pessoas com espectro autista?
\section{Objetivos}
    \subsection{Objetivo Geral}
    Propor um software que auxilia psicólogos na implementação da metodologia ABA para aprendizado por reforço com pacientes do espectro autista.

\subsection{Objetivos Específicos}
\begin{itemize}
    \item Realizar levantamento bibliográfico acerca do uso da metodologia ABA no processo de aprendizagem por reforço;

    \item Identificar e validar funcionalidades relevantes para a implementação da metodologia ABA, com base na análise de trabalhos relacionados e na consulta a profissionais da área;

    \item Analisar a importância de softwares com elevada usabilidade para apoiar a aplicação do método ABA, considerando aspectos de acessibilidade e facilidade de uso;

    \item Desenvolver uma plataforma computacional para apoio ao registro e acompanhamento de dados relacionados à aplicação do método ABA;

    \item Aplicar testes de usabilidade com profissionais da área, a fim de avaliar a interface e as funcionalidades da plataforma.
    % \item Validar a funcionalidade e eficácia da ferramenta por meio de testes de entrada e saída, bem como de testes unitários aplicados nos métodos e funções, a fim de garantir a performance e precisão da aplicação, especialmente em relação aos princípios do Método ABA.
\end{itemize}


